\documentclass[11pt]{article}
\usepackage[T1]{fontenc}
\usepackage{textcomp}
\usepackage[margin=0.75in]{geometry}
\usepackage{amsmath,amssymb,amsthm}
\usepackage{array,booktabs}
\usepackage{hyperref}
\setlength{\parindent}{0pt}
\newtheorem{theorem}{Theorem}
\theoremstyle{definition}
\newtheorem{definition}{Definition}
\title{Flow Proof Artifact: prop-04-side-angle-side}
\author{Generated by \texttt{flow doc proof}}
\date{Flow Proof Book}
\begin{document}
\maketitle
If two triangles have two sides and the included angle equal, the triangles are congruent.\\[0.5em]
\textbf{Source.} euclid — Elements, Book I, Proposition 4\\[0.5em]
\bigskip
\noindent{\large \textbf{Axiom 1.} \textit{If two triangles have two sides and the included angle equal, the triangles are congruent} \hfill the Euclidean plane \cdot Euclid Book I \cdot Proposition 4: side-angle-side congruence}\par
\smallskip\noindent\textit{Coordinate: the Euclidean plane \cdot Euclid Book I \cdot Proposition 4: side-angle-side congruence}\par
\smallskip\noindent\textit{Source: euclid — Elements, Book I, Proposition 4}\par
\medskip
\noindent\textbf{Goal.} If two triangles have two sides and the included angle equal, the triangles are congruent\par
\noindent$\displaystyle \triangle ABC = \triangle DEF$\par
\medskip
\noindent
\begin{tabular}{@{} >{\raggedright\arraybackslash}p{0.06\textwidth} >{\raggedright\arraybackslash}p{0.47\textwidth} >{\raggedleft\arraybackslash}p{0.41\textwidth} @{}}
\textbf{\#} & \textbf{Proof} & \textbf{Mathematics} \\
\midrule
\textbf{1.} & We accept proposition 4: side-angle-side congruence for Book I of the Elements in the Euclidean plane without proof — an ontological commitment, not a lemma. & \\[0.45em]
\textbf{2.} & We invoke the derived fact: segment\_AB == segment\_DE. & \\[0.45em]
\textbf{3.} & We invoke the derived fact: segment\_AC == segment\_DF. & \\[0.45em]
\textbf{4.} & We invoke the derived fact: angle\_BAC == angle\_EDF. & \\[0.45em]
\textbf{5.} & We invoke the derived fact: Common Notion 4: things coinciding with one another are equal. & \\[0.45em]
\textbf{6.} & From step 2, step 3, step 4, and step 5, we can deduce that segment BC equals segment EF. & $\displaystyle \text{segment BC} = \text{segment EF}$ \\[0.45em]
\textbf{7.} & This holds immediately by the stated axiom: angle ABC equals angle DEF. & $\displaystyle \angle ABC = \angle DEF$ \\[0.45em]
\textbf{8.} & This holds immediately by the stated axiom: angle ACB equals angle DFE. & $\displaystyle \angle ACB = \angle DFE$ \\[0.45em]
\textbf{9.} & This holds immediately by the stated axiom: triangle ABC equals triangle DEF. Hence proven. & $\displaystyle \triangle ABC = \triangle DEF$ \\[0.45em]
\bottomrule
\end{tabular}
\label{thm:Geometry-Euclid Book I-proposition_4_side_angle_side_congruence}
\par\medskip\hrule
\end{document}
